% $Id: template.tex 11 2007-04-03 22:25:53Z jpeltier $

\documentclass{vgtc}                          % final (conference style)
% \documentclass[review]{vgtc}                 % review
%\documentclass[widereview]{vgtc}             % wide-spaced review
%\documentclass[preprint]{vgtc}               % preprint
%\documentclass[electronic]{vgtc}             % electronic version

%% Uncomment one of the lines above depending on where your paper is
%% in the conference process. ``review'' and ``widereview'' are for review
%% submission, ``preprint'' is for pre-publication, and the final version
%% doesn't use a specific qualifier. Further, ``electronic'' includes
%% hyperreferences for more convenient online viewing.

%% Please use one of the ``review'' options in combination with the
%% assigned online id (see below) ONLY if your paper uses a double blind
%% review process. Some conferences, like IEEE Vis and InfoVis, have NOT
%% in the past.

%% Figures should be in CMYK or Grey scale format, otherwise, colour 
%% shifting may occur during the printing process.

%% it is recomended to use ``\cref{sec:bla}'' instead of ``Fig.~\ref{sec:bla}''
\graphicspath{{figures/}{pictures/}{images/}{./}} % where to search for the images

\usepackage{times}                     % we use Times as the main font
\renewcommand*\ttdefault{txtt}         % a nicer typewriter font

%% Only used in the template examples. You can remove these lines.
\usepackage{tabu}                      % only used for the table example
\usepackage{booktabs}                  % only used for the table example
\usepackage{lipsum}                    % used to generate placeholder text
\usepackage{mwe}                       % used to generate placeholder figures

%% We encourage the use of mathptmx for consistent usage of times font
%% throughout the proceedings. However, if you encounter conflicts
%% with other math-related packages, you may want to disable it.
\usepackage{mathptmx}                  % use matching math font
\usepackage{amsmath}
\usepackage{amssymb}
\usepackage{array}
\usepackage{multirow}
\usepackage{graphicx}
\usepackage{xcolor}
\usepackage{url}
\usepackage{enumitem}

\newcommand{\R}{\mathbb{R}}
\newcommand{\Domain}{\mathcal{M}}
\newcommand{\Range}{\mathcal{R}}
\newcommand{\Sheets}{\mathcal{S}}
\newcommand{\Sheet}{S}
\newcommand{\Score}{C}
\newcommand{\placeholder}[2]{%
  \begin{center}
  \fbox{%
    \begin{minipage}[c][#1][c]{0.92\linewidth}
      \centering
      #2
    \end{minipage}}
  \end{center}
}
\newcommand{\wideplaceholder}[2]{%
  \begin{center}
  \fbox{%
    \begin{minipage}[c][#1][c]{0.96\textwidth}
      \centering
      #2
    \end{minipage}}
  \end{center}
}
\newcommand{\myparagraph}[1]{\vspace{1mm} \noindent \textbf{#1}}
\newcommand{\ie}{i.e.,\xspace}
\newcommand{\etal}{et al.\xspace}
\newcommand{\Mohit}[1]{\textcolor{blue}{#1}\xspace}
\newcommand{\Talha}[1]{\textcolor{magenta}{Talha: #1}\xspace}
\newcommand{\Nanna}[1]{\textcolor{purple}{#1}\xspace}
\newcommand{\Ingrid}[1]{\textcolor{red}{#1}\xspace}
\newcommand{\Petar}[1]{\textcolor{cyan}{#1}\xspace}

%% If you are submitting a paper to a conference for review with a double
%% blind reviewing process, please replace the value ``0'' below with your
%% OnlineID. Otherwise, you may safely leave it at ``0''.
\onlineid{0}

%% declare the category of your paper, only shown in review mode
\vgtccategory{Research}

%% allow for this line if you want the electronic option to work properly
\vgtcinsertpkg

%% In preprint mode you may define your own headline. If not, the default IEEE copyright message will appear in preprint mode.
%\preprinttext{To appear in an IEEE VGTC sponsored conference.}

%% This adds a link to the version of the paper on IEEEXplore
%% Uncomment this line when you produce a preprint version of the article 
%% after the article receives a DOI for the paper from IEEE
%\ieeedoi{xx.xxxx/TVCG.201x.xxxxxxx}


%% Paper title.

\title{Topological Tracking of Bivariate Features Using Reeb Spaces}

%% This is how authors are specified in the conference style

%% Author and Affiliation (single author).
%%\author{Roy G. Biv\thanks{e-mail: roy.g.biv@aol.com}}
%%\affiliation{\scriptsize Allied Widgets Research}

%% Author and Affiliation (multiple authors with single affiliations).
%%\author{Roy G. Biv\thanks{e-mail: roy.g.biv@aol.com} %
%%\and Ed Grimley\thanks{e-mail:ed.grimley@aol.com} %
%%\and Martha Stewart\thanks{e-mail:martha.stewart@marthastewart.com}}
%%\affiliation{\scriptsize Martha Stewart Enterprises \\ Microsoft Research}

%% Author and Affiliation (multiple authors with multiple affiliations)
\author{Mohit Sharma\thanks{e-mail: mohit.sharma@liu.se}\\ %
        \scriptsize Link\"oping university %
\and Petar Hristov\thanks{e-mail: petar.georgiev.hristov@liu.se}\\ %
        \scriptsize Link\"oping university %
\and Talha Bin Masood\thanks{e-mail: talha.bin.masood@liu.se}\\ %
        \scriptsize Link\"oping university %
\and Ingrid Hotz\thanks{e-mail: ingrid.hotz@liu.se}\\ %
        \scriptsize Link\"oping university %
}

%% A teaser figure can be included as follows
\teaser{
  \centering
  \includegraphics[width=\linewidth]{CypressView}
  \caption{In the Clouds: Vancouver from Cypress Mountain.}
  \label{fig:teaser}
}

%% Abstract section.
\abstract{
Time-varying bivariate fields arise in many scientific applications where the relationship between two scalar quantities changes over time. However, feature-level temporal analysis of such fields remains challenging because bivariate features are not captured by scalar-field extrema, level sets, or contour trees. Reeb spaces provide a topological representation of bivariate fields, but using them for temporal analysis is nontrivial: computed sheets can be noisy, their range-space projections may overlap, and meaningful correspondences across timesteps are difficult to define. Existing continuous-scatterplot-based approaches provide useful timestep-level summaries, including descriptors of domain-segmentation-driven peeled CSP layers, but they do not directly provide a topological feature graph for tracking individual bivariate features. We propose a Reeb-space-based framework for tracking bivariate features over time. At each timestep, we compute the Reeb space, extract prominent sheets by area, and estimate correspondences between sheets in adjacent timesteps using complementary similarity measures in the domain and range. We summarize the resulting temporal correspondences using a Sankey-based sheet evolution summary, linked to sheet renderings, fiber-surface renderings, and diagnostic views. We evaluate the method on a synthetic torus dataset and two time-varying molecular electronic-structure datasets. The results show that Reeb-space sheet tracking reveals persistent features, abrupt reconfiguration intervals, and disagreements between domain- and range-based similarities that are difficult to identify from temporal descriptors of peeled CSP layers alone. Overall, the results demonstrate that Reeb-space sheets can serve as trackable topological features for visual analysis of time-varying bivariate data.


} % end of abstract

%% Keywords that describe your work. Will show as 'Index Terms' in journal
%% please capitalize first letter and insert punctuation after last keyword.
\keywords{Bivariate field, Reeb space, fiber surface, time-varying data, molecular visualization.}

%% Copyright space is enabled by default as required by guidelines.
%% It is disabled by the 'review' option or via the following command:
% \nocopyrightspace

%%%%%%%%%%%%%%%%%%%%%%%%%%%%%%%%%%%%%%%%%%%%%%%%%%%%%%%%%%%%%%%%
%%%%%%%%%%%%%%%%%%%%%% START OF THE PAPER %%%%%%%%%%%%%%%%%%%%%%
%%%%%%%%%%%%%%%%%%%%%%%%%%%%%%%%%%%%%%%%%%%%%%%%%%%%%%%%%%%%%%%%%

\begin{document}

%% The ``\maketitle'' command must be the first command after the
%% ``\begin{document}'' command. It prepares and prints the title block.

%% the only exception to this rule is the \firstsection command
\firstsection{Introduction}

\maketitle
Many scientific simulations produce multiple scalar fields over a common spatial domain. In molecular electronic-structure analysis, for example, an electronic transition can be represented using hole and particle natural transition orbital (NTO) fields~\cite{Martin2003NTO}. These fields indicate where charge is depleted and accumulated during the transition. As the molecular geometry evolves, the analysis becomes a time-varying bivariate problem: the goal is not only to study a pair of fields, but also to understand how their joint features evolve over time.

A common strategy for analyzing bivariate fields is to inspect selected isosurfaces or fiber surfaces. Fiber surfaces generalize isosurfaces to bivariate data and provide detailed local views of structures associated with chosen range values. They are useful for examining specific hypotheses, but they do not by themselves provide a compact representation of temporal feature evolution. Another strategy is to summarize each timestep using a continuous scatterplot (CSP). CSPs generalize discrete scatterplots to continuous spatial fields and summarize the distribution of bivariate values. Recent work has used CSP descriptors to analyze the temporal evolution of molecular electronic-structure data~\cite{Sharma2025}. In that approach, bivariate features are obtained using domain-segmentation-driven CSP peeling~\cite{Sharma2021}. The resulting descriptors capture the temporal behavior of the peeled CSP layers. However, they do not directly provide a topological feature evolution that identifies how individual bivariate features persist, split, disappear, merge, or reconfigure across timesteps, nor do they show which domain vertices continue to participate in the same evolving bivariate feature.

We address this limitation using Reeb spaces. For a bivariate field $F=(f,g):\Domain\rightarrow\R^2$, the Reeb space is obtained by identifying domain points that lie in the same connected component of a fiber $F^{-1}(a,b)$. Its two-dimensional cells, called sheets, combine range-space geometry with domain connectedness. A sheet can therefore be interpreted as a coherent bivariate feature. Tracking sheets across timesteps provides a feature-level representation of temporal evolution.


This paper presents a practical visual analytics pipeline for Reeb-space sheet tracking. At each timestep, we compute a Reeb space and retain the top sheets by range-space area. We compare sheets between adjacent timesteps using range-space shape similarity and domain-space vertex overlap. The resulting temporal graph is explored through a Sankey-based view, linked sheet images, linked fiber-surface images, and diagnostic plots. The tool does not force a one-to-one matching, because Reeb-space sheets can split, merge, appear, disappear, and reconfigure. Instead, it provides candidate correspondences and lets the user compare how different continuity notions behave.

The key idea is that range-space and domain-space continuity answer different questions. Range-space similarity asks whether a sheet preserves a similar relationship between the two scalar fields. Domain-space overlap asks whether the same domain vertices continue to participate in the sheet. These two notions of continuity need not agree. We therefore report range-space events and domain-space events separately, each ranked by its corresponding event score. In addition, we separately report domain--range target disagreements. This disagreement analysis does not replace the range or domain event rankings. Instead, it provides a way to select concrete cases where the two continuity notions confidently identify different target sheets, highlighting intervals where range-space behavior and domain participation diverge. The main contributions of this paper are:
\begin{itemize}[leftmargin=*]
  \item A Reeb-space sheet-tracking workflow for time-varying bivariate fields, where prominent sheets are treated as trackable topological features.

  \item Complementary correspondence measures for range-space sheet geometry and domain-vertex overlap, together with diagnostic summaries that preserve their distinct interpretations.

  \item A linked visual interface that combines Sankey-based temporal summaries, event-score plots for highlighting potentially interesting timesteps, continuing-feature lists, range-space sheet renderings, fiber-surface renderings, and domain--range disagreement highlights.

  \item Case studies on two molecular datasets, MVK and stilbene, and one synthetic torus dataset, showing that range-space sheet events reveal meaningful reconfiguration intervals, while domain-space events provide a complementary view of changes in domain-vertex participation.
\end{itemize}





\begin{figure*}[t]
    \centering
    \includegraphics[width=0.9\textwidth]{./figures/rs-background.pdf}
    \caption{Introduction to the main topological and geometric structures used in the paper.
    We assume that our domain \emph{a)} is a tetrahedral mesh and we render a fiber and a fiber surface.
    In the range \emph{b)} we show the arrangement of the images of the edges of the tetrahedra. 
    Combining the faces of this arrangement yields the three sheets of the Reeb space \emph{c)}.
    Finally, the sheets of the Reeb space correspond to three volumes which segment the domain.
    We call those volumes features.
    }
    \label{fig:pipeline}
\end{figure*}

\section{Related work}
\Talha{First draft largely taken from our previous papers. TODO: Update and shorten; Add a few more up-to-date references, and fix issues with bibtex entries.}

\myparagraph{Reeb space.} The Reeb space is a higher-dimensional generalization of the Reeb graph that captures the topology of multiple scalar functions defined on a common domain \cite{Edelsbrunner2008}. Whereas a Reeb graph summarizes the connectivity of level sets of a single scalar field, the Reeb space describes the connectivity of \emph{fibers} in multivariate data. In the bivariate case, it consists of a collection of two-dimensional sheets connected along Jacobi edges \cite{Edelsbrunner2004b}. This structure provides a compact representation of the topological organization of a multifield and forms the basis of many approaches to multivariate topological analysis. Several algorithms have been proposed for computing Reeb spaces in bivariate settings \cite{Tierny2017,chattopadhyayAlgorithmFastCorrect2024,Hristov2026SingularAA}, as well as for more general domains and higher-dimensional ranges \cite{Edelsbrunner2008,hristov2025}. One of the key challenges in these methods is achieving robustness in the presence of numerical inaccuracies \cite{Hristov2025b}. To improve scalability, approximate representations such as the Joint Contour Net \cite{Carr2014} and Mapper \cite{Dey2016,Brown2021} have also been developed, providing some practical alternatives for the analysis of complex multivariate datasets. In addition to serving as topological descriptors, Reeb spaces and their approximations have been used to compare multifields by defining distances between their topological structures \cite{Ramamurthi2022a,Agarwal2021,Ramamurthi2024}.

\myparagraph{Approaches to feature tracking.} Topology-based feature tracking is a central task in topological data analysis and visualization of scalar and vector fields. Broadly, tracking approaches follow two paradigms: either features are extracted independently at each time step and subsequently matched, or the temporal dimension is incorporated directly into the data domain, transforming tracking into a spatio-temporal feature extraction problem \cite{Post2003,Weber2011}. Topological features are commonly represented either by critical points, such as maxima and minima, or by spatial regions defined through sub-, super-, or level sets. Intermediate feature definitions also exist, including groups of critical points represented by merge-tree structures, such as crown features, which have been employed in applications ranging from cyclone tracking to medical imaging \cite{Engelke2020,Nilsson2022Cyclone,Rasheed2022}. Other application-specific descriptors include vortex features derived from scalar quantities in flow fields and extremal structures based on subsets of the Morse-Smale complex.

A large body of work focuses on tracking critical point-based features. Early approaches tracked vortices and other flow phenomena through critical points in derived scalar fields \cite{Garth2004b,Weinkauf2011}. Reininghaus et al.~\cite{Reininghaus2011} introduced Combinatorial Feature Flow Fields, a history-based framework that has since been extended to track hierarchical critical point structures and applied in several domains \cite{Engelke2020,Nilsson2022Cyclone}. More recently, correspondence has been established through topological similarity measures, including persistence-diagram matching, merge-tree mappings, and tree-edit distances \cite{Cohen-Steiner2010,Sridharamurthy2020,Wetzels2022a,Wetzels2023}. Alternative approaches combine topological and geometric information, exploit temporal interpolation, or use optical-flow-based correspondences \cite{Yan2022b,Valsangkar2019}. Critical points have also been tracked indirectly through their associated ascending and descending manifolds in the Morse-Smale complex \cite{Nilsson2023}.

For spatial features, correspondence is typically determined through similarity measures such as spatial distance, attribute similarity, and, most commonly, spatial overlap. Foundational work by Samtaney et al.~\cite{Samtaney1994} and Silver and Wang~\cite{silver1998tracking} established overlap-based tracking of threshold-defined regions, augmented by geometric attributes such as volume and center of mass. Spatial overlap remains widely used, including in methods based on merge trees, where hierarchical feature relationships can be exploited for tracking and visualization through nested tracking graphs \cite{Widanagamaachchi2012,Saikia2017,Lukasczyk2017nested}. Additional processing steps, such as semantic filtering or global optimization, are often employed to remove implausible correspondences and construct coherent feature tracks \cite{Nilsson2022Cyclone,Schnorr2020}. Hierarchical tracking methods leverage the nested structure of level sets and merge trees to capture multi-scale behavior \cite{Lukasczyk2019dynamic}, while layout and visualization techniques have improved the readability of the resulting tracking graphs \cite{Kopp2019,Dobler2024}. In contrast to traditional applications that focus on sparse features occupying a small fraction of the domain, Schnorr et al.~\cite{Schnorr2020} considered space-filling features in combustion simulations, where the entire domain is partitioned into regions. Although overlap-based correspondence remains effective in this setting, it produces highly connected tracking graphs that require computationally expensive optimization procedures to extract meaningful feature trajectories.


\section{Background}
This section introduces the concepts used throughout the paper. We begin with bivariate fields and their domain and range spaces, then review fibers, fiber surfaces, Reeb spaces, and Reeb space sheets. These concepts provide the basis for formulating sheet tracking as a temporal correspondence problem between prominent Reeb space sheets and their corresponding features.

\myparagraph{Bivariate fields.} Let $\Domain$ be a triangulation of a regular-grid volume or more generally a 3-manifold.
A \textit{scalar field} over $\Domain$ is a piecewise-linear (PL) map $f:\Domain \to \R$, typically defined on the vertices of $\Domain$ and barycentricaly interpolated elsewhere.
A \textit{bivariate field} is a PL map $F=(f,g):\Domain \to \R^2,$
where $f$ and $g$ are scalar fields defined on the same domain.
We refer to $\Domain$ as the \textit{domain space}, to $\R^2$ as the \textit{range space} and to $F(\Domain)$ as the \textit{image} of $\Domain$ under $F$.
A bivariate field is \textit{generic} when the image of no three vertices of $\Domain$ are collinear in the plane and the images of no three edges of $\Domain$ are concurrent in their interiors~\cite{Hristov2025b}.
% This ensures that preimages are well-behaved.

\myparagraph{Fibers.} The preimage of a point $y=(a,b)\in\R^2$ defines a \emph{fiber}
$F^{-1}(y)=\{x \in \Domain \mid f(x)=a,\ g(x)=b\}$.
The point $p$ is referred to as a \emph{control point}.
Generally, the fibers of $F$ are collections of one-dimensional curves.
Fibers can only change their connectivity at singular edges and vertices~\cite{hristovHypersweepsConvectiveClouds2022}, which generalize critical vertices in the scalar case.
The preimage $F^{-1}(P)$ of a polygon or line $P \subset \R^2$ in the range is a collection of fibers that sweep out a \emph{fiber surface}~\cite{Carr2015}.
We refer to $P$ as a \emph{control polygon or line}.
Fiber surfaces generalize isosurfaces in that they are used to define boundaries of features when visualizing bivariate data.


\myparagraph{Reeb spaces.} The Reeb space generalizes the Reeb graph from scalar functions to multivariate maps~\cite{Edelsbrunner2008}.
It captures how the connected components of fibers change for all points over the range of $F$.
The Reeb space is the quotient space $\mathcal{R}_F = \Domain/\sim$, where $\sim$ is the equivalence relation on $\Domain$ such that $x\sim x' \text{ if } F(x) = F(x') \text{ and }
x, x' \text{ lie in the same connected component of } F^{-1}(F(x))$.
Practical algorithms for computing the Reeb space include Jacobi fiber surfaces~\cite{Tierny2017b} and the arrange and traverse~\cite{hristov2025, Hristov2026SingularAA} framework.

\myparagraph{Reeb Space Sheets.} 
For a bivariate map over a three-dimensional domain, the Reeb space is represented by a two-dimensional cell complex, whose 2-cells are referred to as \textit{sheets}.
The sheets of the Reeb space are glued along singular vertices and edges where the connectivity of fibers changes, called the Jacobi structure~\cite{Chattopadhyay2014}.
Preimages of sheets in the domain correspond to features captured by the Reeb space, analogous to how to preimages of edges of the Reeb graph and contour tree correspond to features in scalar fields.

\myparagraph{Sheet and Feature Metrics.}
Each sheets of the Reeb space can be described geometrically by its projection onto the range space as a polygon.
In the arrange and traverse framework~\cite{hristov2025, Hristov2026SingularAA}, each such polygon is the union of faces of the geometric arrangement of the images of singular edges.
The area of each sheet is then the sum of areas of its associated arrangement faces.
The volume of a feature associated with a sheet can be computed exactly using the Jacobi fiber surfaces that bound it or approximated with the number of regular points in that volume (similarly to the scalar case~\cite{Schneider2008}).

\section{Method}
This section describes our pipeline for tracking Reeb space sheets in time-varying bivariate data. Starting from a sequence of volumetric datasets with two scalar fields, we compute the Reeb space at each timestep and extract the most prominent sheets. We then generate visual representations of their range-space geometry and corresponding spatial fiber-surface context. To track sheets over time, we compare sheets from adjacent timesteps using complementary similarity measures in the range and domain. The resulting correspondences define a temporal graph, which we visualize and analyze to identify persistent sheets, weak continuations, and candidate transition intervals.

\subsection{Input and Preprocessing}

\begin{figure*}[t]
  \wideplaceholder{1.55in}{Figure placeholder: pipeline overview. Show input
  time-varying bivariate fields, Reeb space computation, top-sheet extraction,
  sheet and fiber-surface rendering, pairwise sheet matching, event analysis,
  and Sankey-style visual exploration.}
  \caption{Overview of the proposed Reeb space sheet tracking pipeline.}
  \label{fig:pipeline}
\end{figure*}

The input to our method is a time-ordered sequence of volumetric datasets. Each timestep contains two scalar fields defined on a common spatial domain, forming a bivariate field $F_i=(f_i,g_i)$. We assume that the field pair is consistently defined across the sequence so that the corresponding scalar quantities are comparable over time.
\Mohit{Mention how is sheet area computed?}
For each timestep, we compute the Reeb space of the bivariate field and extract its sheet representation. For each sheet, this representation stores its geometric footprint in the range, its range-space area, and the associated set of regular domain vertices, when available. Some sheets may have no associated regular vertices due to input sampling, or the structure of the computed sheet representation. We retain such sheets because they may still represent meaningful range-space geometry. The extracted sheet information provides the basis for sheet ranking, visualization, and correspondence estimation.




Only timesteps for which a valid Reeb space sheet representation is obtained are included in the tracking stage. Timesteps where the Reeb space computation fails, for example due to high degeneracy, are excluded from the temporal correspondence graph.



\subsection{Top-Sheet Selection}

Let $\Sheets_i^\mathrm{all}$ denote the set of all sheets extracted at timestep $i$. We rank these sheets by their range-space area and retain the $N$ largest sheets:
\[
\Sheets_i = \operatorname{TopN}(\Sheets_i^\mathrm{all}, N)
\]
This selection focuses the analysis on dominant bivariate features while keeping the temporal graph manageable. The value of $N$ is a user-defined parameter and can be varied to study the sensitivity of the analysis to smaller or less prominent sheets.


\begin{figure}[t]
  \placeholder{1.35in}{Figure placeholder: one timestep with the full
  Reeb space sheet collection, top sheets highlighted, and a zoomed example of
  one sheet footprint in range space.}
  \caption{A Reeb space sheet is represented by its footprint in the bivariate
  range and by its associated domain support.}
  \label{fig:sheet-definition}
\end{figure}


\subsection{Sheet Correspondence Measures}

Given a time-varying sequence of bivariate fields $F_0,F_1,\ldots,F_T,$
let $\Sheets_i$ denote the set of sheets extracted from the Reeb space of $F_i$. The goal is to identify how sheets evolve between consecutive timesteps by estimating candidate correspondences between sheets in $\Sheets_i$ and $\Sheets_{i+1}$. For each adjacent timestep pair $(i,i+1)$, we compare every source sheet $\Sheet\in\Sheets_i$ with every target sheet $\Sheet'\in\Sheets_{i+1}$. A correspondence indicates that two sheets may represent the same evolving bivariate feature, or related parts of a feature before or after a topological change.

We store all candidate correspondences rather than solving a global one-to-one assignment. This design is intentional because Reeb-space sheets may persist, split, merge, appear, disappear, or change substantially in shape and spatial overlap. Therefore, for each adjacent timestep pair, we construct a weighted bipartite graph whose nodes are sheets from $\Sheets_i$ and $\Sheets_{i+1}$, and whose edge weights encode sheet similarity. Retaining the full candidate graph supports visual inspection of continuations, splits, merges, and uncertain correspondences. We estimate these edge weights using complementary similarity measures defined in the range and in the domain.

\myparagraph{Range-space similarity.}
Range-space similarity measures how similar two sheets are as geometric objects in the bivariate range. The primary range-space measure is shape intersection-over-union (iou). Let $M_\Sheet$ and $M_{\Sheet'}$ denote binary masks representing the range-space footprints of two sheets under a common coordinate system. We define
$$
q_\mathrm{iou}(\Sheet,\Sheet') =
\frac{|M_\Sheet \cap M_{\Sheet'}|}
{|M_\Sheet \cup M_{\Sheet'}|}
$$
This score compares the actual sheet footprints in range space and therefore captures overlap, relative position, relative size, and detailed shape under a common coordinate frame. We use $q_\mathrm{iou}$ as the default range-space correspondence score $q_R$ in our analysis.

The implementation also computes auxiliary range-space scores, including area ratio, bounding-box IoU, and centroid similarity. These scores are useful as diagnostics and optional viewer metrics, but we do not use them as primary result scores. Area ratio ignores position, centroid similarity ignores sheet shape, and bounding-box IoU is coarser than the actual sheet footprint. Their definitions and agreement statistics are provided in the supplement.

\myparagraph{Domain-space similarity.}
Domain-space similarity compares the spatial regions associated with sheets across adjacent timesteps. This measure complements range-space similarity: two sheets may occupy similar regions in the bivariate range while being supported by different domain vertices, and conversely, sheets with overlapping domain regions may change their range-space geometry.

Let $V_\Sheet$ and $V_{\Sheet'}$ denote the sets of regular domain vertices associated with a source sheet $\Sheet$ and a target sheet $\Sheet'$, respectively. We first compute their raw spatial overlap: $ O(\Sheet,\Sheet') = |V_\Sheet \cap V_{\Sheet'}|$.
A positive overlap indicates that the two sheets share part of their spatial region and are therefore candidates for temporal correspondence. To account for changes in sheet size, we compute two directional normalized overlap scores:
$$
d_\mathrm{src}(\Sheet,\Sheet') =
\begin{cases}
O(\Sheet,\Sheet')/|V_\Sheet|, & |V_\Sheet|>0,\\
0, & |V_\Sheet|=0
\end{cases}
$$
and
$$
d_\mathrm{tgt}(\Sheet,\Sheet') =
\begin{cases}
O(\Sheet,\Sheet')/|V_{\Sheet'}|, & |V_{\Sheet'}|>0,\\
0, & |V_{\Sheet'}|=0
\end{cases}
$$
The source-normalized score $d_\mathrm{src}$ measures how much of the source sheet's spatial region continues into the target sheet, while the target-normalized score $d_\mathrm{tgt}$ measures how much of the target sheet's spatial region is inherited from the source sheet. We summarize the directional overlap scores using
$$
q_D(\Sheet,\Sheet') =
d_\mathrm{max}(\Sheet,\Sheet') =
\max\left(
d_\mathrm{src}(\Sheet,\Sheet'),
d_\mathrm{tgt}(\Sheet,\Sheet')
\right)
$$
This score lies in $[0,1]$ and treats containment as strong evidence of domain-space continuity. For example, if a small source sheet is fully contained in a larger target sheet, then $d_\mathrm{src}$ is high even though the target sheet contains additional vertices. This behavior is useful for identifying continuations involving growth, shrinkage, merging, or absorption. A natural symmetric alternative is the vertex-set Jaccard overlap,
$$
q_\mathrm{jac}(\Sheet,\Sheet') =
\frac{|V_\Sheet\cap V_{\Sheet'}|}
{|V_\Sheet\cup V_{\Sheet'}|}
$$
However, Jaccard overlap penalizes size imbalance between sheets. The Jaccard score would be low when a larger target sheet contains many vertices outside a smaller source sheet, even if the source sheet is fully contained within it.  This is undesirable in our setting because containment relationships can indicate meaningful temporal behavior. We therefore use $d_\mathrm{max}$ as the default domain-space correspondence score $q_D$.

In the visual interface, link width and opacity can be mapped to different domain-space quantities, including $O(\Sheet,\Sheet')$, $d_\mathrm{src}$, $d_\mathrm{tgt}$, and $d_\mathrm{max}$. This allows the analyst to inspect raw overlap as well as directional overlap. However, all domain-space analyses reported in the paper use $d_\mathrm{max}$ unless otherwise stated.

\subsection{Diagnostic Summaries}

The sheet correspondence graph can contain many candidate links, especially for long time sequences. To guide inspection, we compute three diagnostic summaries: event scores for timestep intervals, lifetimes for continuing features, and domain--range target disagreement scores. These diagnostics are used for ranking and visual exploration. They are not intended as formal topological event detectors.

\myparagraph{Event intervals.}
Let $q$ denote the active correspondence score, either the range-space score $q_R$ or the domain-space score $q_D$. For an adjacent timestep pair $(i,i+1)$, let $N_i=|\Sheets_i|$. For a source sheet $\Sheet\in\Sheets_i$, we define its best outgoing score as
$$
B^\mathrm{out}(\Sheet)=
\max_{\Sheet'\in\Sheets_{i+1}} q(\Sheet,\Sheet')
$$
This is the strongest available continuation of $\Sheet$ into the next timestep. Similarly, for a target sheet $\Sheet'\in\Sheets_{i+1}$, we define its best incoming score as
$$
B^\mathrm{in}(\Sheet')=
\max_{\Sheet\in\Sheets_i} q(\Sheet,\Sheet')
$$
This is the strongest available predecessor of $\Sheet'$ from the previous timestep.

Given a threshold $\theta$, we count source sheets whose best continuation is weak:
$$
W^\mathrm{out}_i(\theta)=
\left|
\left\{
\Sheet\in\Sheets_i:
B^\mathrm{out}(\Sheet)<\theta
\right\}
\right|
$$
We also count target sheets whose best predecessor is weak:
$$
W^\mathrm{in}*i(\theta)=
\left|
\left\{
\Sheet'\in\Sheets*{i+1}:
B^\mathrm{in}(\Sheet')<\theta
\right\}
\right|
$$
Large values of these two terms indicate intervals where many sheets lack strong continuations or predecessors.

We next count possible branching configurations. A source sheet is counted as a possible split if it has at least two target sheets with above-threshold scores:
$$
P^\mathrm{split}_i(\theta)=
\left|
\left\{
\Sheet\in\Sheets_i:
\left|
\left\{
\Sheet'\in\Sheets_{i+1}:
q(\Sheet,\Sheet')\geq\theta
\right\}
\right|
\geq 2
\right\}
\right|
$$
Similarly, a target sheet is counted as a possible merge if it has at least two source sheets with above-threshold scores:
$$
P^\mathrm{merge}_i(\theta)=
\left|
\left\{
\Sheet'\in\Sheets_{i+1}:
\left|
\left\{
\Sheet\in\Sheets_i:
q(\Sheet,\Sheet')\geq\theta
\right\}
\right|
\geq 2
\right\}
\right|
$$
These terms capture correspondence ambiguity. Such ambiguity may indicate split or merge behavior, but it may also arise when several sheets have similar range-space geometry or overlapping domain support.

Finally, we compute the mean best outgoing score:
$$
\overline{B^\mathrm{out}_i}=
\frac{1}{N_i}
\sum_{\Sheet\in\Sheets_i}
B^\mathrm{out}(\Sheet)
$$
The event score for interval $i\rightarrow i+1$ is defined as
$$
E_i(\theta)=
W^\mathrm{out}_i(\theta)
+
W^\mathrm{in}_i(\theta)
+
\frac{1}{2}P^\mathrm{split}_i(\theta)
+
\frac{1}{2}P^\mathrm{merge}_i(\theta)
+
\left(1-\overline{B^\mathrm{out}_i}\right)N_i
$$
This is a ranking score, not a probability. The weak outgoing and weak incoming terms emphasize missing or weak continuations. The split and merge terms emphasize branching ambiguity, but with half weight because such sheets still have above-threshold correspondences. The final continuation-gap term increases the score when the average best continuation is low, even if some individual sheets remain above threshold. Thus, a large value of $E_i(\theta)$ indicates an interval where temporal continuity is weak, ambiguous, or both.

Range-space event intervals use $q=q_R$ and answer whether the range-space sheet footprints change. Domain-space event intervals use $q=q_D$ and answer whether the overlap among associated domain vertices changes. These two rankings are reported separately because they capture different notions of continuity. For range IoU, we use $\theta=0.5$ as the main threshold, since it corresponds to a substantial footprint overlap under a common range-space frame. For domain overlap, the interpretation of a fixed threshold is more dataset-dependent because vertex-set sizes, mesh sampling, and containment behavior can vary. Therefore, the viewer also supports percentile-based domain thresholds, such as $Q_{25}$, $Q_{50}$, $Q_{75}$, and $Q_{90}$, computed from the best available domain continuations. The he
default domain threshold is the median $Q_{50}$. These thresholds allow the analyst to inspect typical, strong, and conservative domain continuities. Sensitivity summaries are exported for multiple $\theta$ vlaues for both domain and range.

\myparagraph{Continuing-feature lifetimes.}
To identify persistent sheet behavior, we compute a greedy lifetime diagnostic. For a sheet $\Sheet_i\in\Sheets_i$, we select the best target sheet in the next timestep:
$$
\Sheet_{i+1}^{\ast}=
\arg\max_{\Sheet'\in\Sheets_{i+1}}
q(\Sheet_i,\Sheet')
$$
If the best score satisfies
$$
q(\Sheet_i,\Sheet_{i+1}^{\ast})\geq\theta
$$
the track is extended to $\Sheet_{i+1}^{\ast}$, and the same rule is applied again from the next timestep. The lifetime is the number of timesteps visited before the best continuation falls below the threshold or the sequence ends.

For each greedy track, we record its length, mean continuation score, and minimum continuation score. This diagnostic helps identify long-lived or stable sheets. However, it does not enforce a global one-to-one assignment, and multiple starting sheets may converge to the same later sheet. Therefore, lifetimes should be interpreted as evidence of local temporal persistence, not as definitive physical feature trajectories.

\myparagraph{Domain--range target disagreement.}
Range-space and domain-space continuity can suggest different continuations for the same source sheet. Range-space similarity asks whether the sheet preserves a similar footprint in the bivariate range. Domain-space overlap asks whether the same domain vertices continue to participate in the sheet. These two questions may have different answers: a sheet may keep a similar range-space shape while moving to a different spatial region, or it may keep much of its spatial region while changing its range-space footprint.

We therefore compute a target disagreement diagnostic. The idea is simple: for each source sheet, we ask which target sheet is preferred by the range-space score and which target sheet is preferred by the domain-space score. If both scores choose the same target, the two notions agree. If they choose different targets, we measure how strongly they disagree.

For each source sheet $\Sheet\in\Sheets_i$, let
$$
T_R(\Sheet)=
\arg\max_{\Sheet'\in\Sheets_{i+1}}
q_R(\Sheet,\Sheet')
$$
be the target preferred by range-space similarity, and let
$$
T_D(\Sheet)=
\arg\max_{\Sheet'\in\Sheets_{i+1}}
q_D(\Sheet,\Sheet')
$$
be the target preferred by domain-space overlap. If $T_R(\Sheet)=T_D(\Sheet)$, the two scores agree for that source sheet. If $T_R(\Sheet)\neq T_D(\Sheet)$, the source sheet contributes a disagreement example.

To quantify the disagreement, we compare each score with what it would obtain if it were forced to accept the other score's preferred target. Let
$$
q_R^\ast=q_R(\Sheet,T_R(\Sheet)),
\qquad
q_D^\ast=q_D(\Sheet,T_D(\Sheet))
$$
be the best scores achieved by the range and domain choices. We then evaluate the cross choices:
$$
q_R^D=q_R(\Sheet,T_D(\Sheet)),
\qquad
q_D^R=q_D(\Sheet,T_R(\Sheet))
$$
Here, $q_R^D$ is the range-space score of the domain-preferred target, and $q_D^R$ is the domain-space score of the range-preferred target.

The disagreement is meaningful only if accepting the other target causes a clear loss. We therefore define the range and domain losses as
$$
L_R=\max(0,q_R^\ast-q_R^D),
\qquad
L_D=\max(0,q_D^\ast-q_D^R)
$$
A large $L_R$ means that the range-space score strongly prefers its own target over the domain-selected target. A large $L_D$ means that the domain-space score strongly prefers its own target over the range-selected target.

We also include a confidence term,
$$
C=\min(q_R^\ast,q_D^\ast)
$$
which is high only when both scores have strong preferred targets. This prevents weak or uncertain matches from producing large disagreement values. The source-level disagreement score is
$$
D(\Sheet)=
\frac{1}{2}(L_R+L_D)C
$$
Thus, $D(\Sheet)$ is high only when range-space and domain-space continuity confidently choose different targets and each score loses substantially by accepting the other score's target. It is low when the two targets are nearly interchangeable, or when one of the two scores is already weak.

For each timestep pair, the viewer summarizes these source-level disagreements using the disagreement count, disagreement fraction, mean disagreement score, and maximum disagreement score:
$$
D_\mathrm{max}(i,i+1)=
\max_{\Sheet\in\Sheets_i}D(\Sheet)
$$
The domain--range disagreement plot uses $D_\mathrm{max}$ on the vertical axis. A high value means that at least one source sheet in that interval has a strong local disagreement between range-space and domain-space continuation evidence. It does not mean that all sheets in the interval disagree, so this plot should be interpreted together with the disagreement count, agreement fraction, and the corresponding Sankey links.





\subsection{Sheet and Fiber-Surface Renderings}
For interactive analysis, we precompute visual representations of the selected sheets and their corresponding spatial context. Each selected sheet is represented by its footprint as a polygonal region in the bivariate range. These range-space footprints are used for both visualization and shape-based comparison across timesteps. Each sheet is rendered with a uniform color so that the image emphasizes its geometry. The full Reeb space is shown as faint context, while the selected sheet remains visually dominant.

To support meaningful visual comparison over time, all sheet renderings use common range-space bounds and a fixed image size. This ensures that apparent changes in sheet size and position reflect changes in the data, rather than independent rescaling of each rendered image.

We also precompute fiber-surface renderings for selected sheet and timestep combinations. These renderings provide spatial context by showing how a selected range-space sheet is embedded in the domain. The specific fiber surfaces are chosen according to the dataset and analysis task, and are stored for inspection alongside the corresponding Reeb space sheets. Together, the sheet and fiber-surface renderings help relate temporal sheet correspondences to the underlying volumetric structures.

\subsection{Temporal Graph and Sankey Layout}
The main visualization is a Sankey-style temporal graph. Timesteps are placed from left to right, selected sheets are shown as nodes, and candidate correspondences are shown as links. Link thickness and opacity encode the active correspondence score after normalization within the selected metric mode. Analysts can filter links by threshold, display only the strongest outgoing link per source sheet, change the visible timestep range and stride, adjust the number of top sheets, and reorder sheets using weighted crossing-reduced ordering, sheet area, or domain vertex count.

The interface is designed for comparison and Multiple Sankey panels can be shown side by side.

The analysis panel turns the graph into an exploratory workflow. The interesting interval tab ranks timestep based on weak continuations or possible split/merge ambiguity. The continuing-feature tab follows best outgoing links above the selected threshold ($\theta$). The sensitivity tab repeats these summaries over different $\theta$ values and helps identify events or timesteps that are robust across multiple ($\theta$). The domain/range complementarity tab highlights source sheets for which range and domain criteria confidently choose different targets. Selection in the plots in these anlyses tabs is linked to the snakey diagram and the corresponding timsteps/features get highlighted. Node and link selections open linked detail views: a node shows sheet metadata, the rendered sheet footprint, and available fiber-surface images; a link shows source and target sheets side by side. Image views support linked pan and zoom so that corresponding sheet or fiber-surface structures can be compared at the same scale. The supplement gives the exhaustive control-level description.


\begin{figure*}[t]
  \wideplaceholder{1.75in}{Figure placeholder: unified viewer screenshot.
  Include a Sankey diagram with node/link coloring, a selected link, side-by-side
  sheet images, and side-by-side fiber-surface images with linked zoom.}
  \caption{The visual interface links temporal sheet correspondences with
  range-space sheet images and spatial fiber-surface evidence.}
  \label{fig:viewer}
\end{figure*}

\section{Results}
We evaluate our Reeb-space sheet-tracking framework on one controlled synthetic dataset and two molecular electronic-structure datasets. The synthetic torus dataset contains $100$ timesteps and provides a controlled setting for testing whether the method captures a known transition in sheet structure. The molecular datasets are obtained from photoinduced molecular dynamics simulations and include methylvinylketone (MVK) with $83$ timesteps and \textit{cis}-stilbene with $704$ timesteps. In the viewer, the MVK data is represented by two processed datasets, MVK\_s1 and MVK\_s2, corresponding to the two electronic states.

The molecular case studies follow the setting of previous bivariate analyses of electronic-structure evolution~\cite{Sharma2025}. In these data, pairs of scalar fields represent the hole and particle natural transition orbitals (NTOs)~\cite{martin_nto_2023} over time. The two NTO fields describe where electronic density is depleted and where it is accumulated for a given excited state relative to the ground state. We treat these fields as a time-varying bivariate field over evolving nuclear geometries. The goal is to support feature-level analysis of this evolution. Our analysis focuses on prominent Reeb-space sheets extracted at each timestep, their temporal correspondences, and diagnostic summaries that identify persistent features and intervals of possible reconfiguration.

For all datasets, we compute the top $N=20$ sheets at each timestep. The viewer displays the top $10$ sheets by default for readability, but this can be increased up to $20$. We report range-space events as the primary signal of feature reconfiguration, because the default range correspondence score is based on sheet-footprint IoU in the bivariate range. Domain-space event rankings are reported separately and interpreted as changes in the overlap of associated domain vertices. We also use domain--range disagreement diagnostics to identify concrete examples where range-space and domain-space continuity choose different target sheets. This separation is important because the two scores answer different questions: range-space similarity asks whether the bivariate relationship persists, whereas domain-space overlap asks whether the same spatial region continues to participate in that relationship.

\subsection{Bivariate torus}
\subsection{Methylvinylketone}
\subsection{\textit{cis}-stilbene}





%% if specified like this the section will be committed in review mode
\acknowledgments{}

%\bibliographystyle{abbrv}
\bibliographystyle{abbrv-doi}
%\bibliographystyle{abbrv-doi-narrow}
%\bibliographystyle{abbrv-doi-hyperref}
%\bibliographystyle{abbrv-doi-hyperref-narrow}

\bibliography{bib}
\end{document}
